\section{Structuring lemmas for reuse}\label{sec:13-working-efficiently:05-structuring-lemmas:1}

The single biggest efficiency gain, greater than any tactic choice, is to \textbf{prove the general fact once, as its own named lemma, as soon as an argument would otherwise be repeated.} \texttt{left\_\allowbreak{}inverse\_\allowbreak{}unique} from Chapter 8 is the running example. Theorem 3 (\texttt{inv\_\allowbreak{}op}), \texttt{neg\_\allowbreak{}one\_\allowbreak{}mul}, and the exercise \texttt{neg\_\allowbreak{}mul} from Chapter 10 all reduce to it, instead of re-deriving ``uniqueness of inverses'' inline. Signs that a lemma should be factored out are as follows.

\begin{itemize}
\tightlist
\item
  A \href{https://lean-lang.org/doc/reference/latest/Tactic-Proofs/Tactic-Reference/}{\texttt{rw}} chain already written for a different but structurally identical goal is about to be repeated. In that case, the shared shape should instead be stated as its own \texttt{theorem}/\texttt{have}, then applied via \texttt{apply}/\texttt{exact} in both places.
\item
  A sub-goal deep in a proof would itself be a reasonable, independently statable mathematical fact (for example, ``an element that equals its own double is zero,'' buried inside \texttt{mul\_\allowbreak{}zero} of Chapter 10). Naming it is both more efficient \emph{and} more readable, since the outer proof then reads as a short chain of named facts instead of one long undifferentiated block.
\end{itemize}

This is the same judgment call made writing ordinary code, extract a helper when, and only when, real duplication or a genuinely separate sub-claim is present, not ahead of time ``just in case.''

\textbf{Key points.} \texttt{exact?}/\texttt{apply?} search for a closing term but do not always find the shortest one, a correct but roundabout result is still worth simplifying by hand afterward, and still cheaper to obtain that way than by deriving the whole term from nothing. \texttt{decide}/\texttt{omega}/\texttt{norm\_\allowbreak{}num} replace a hand proof exactly on their decidable fragment: a finite, concretely enumerable carrier such as the \texttt{Fin 3} of \texttt{fin3Group}/\texttt{fin3Ring} (Chapter 9), never a goal about a generic, unspecified structure such as an arbitrary \texttt{Group G} or \texttt{Ring R} (Chapters 8 and 10), where no decision procedure applies and a hand-built proof is the only option. \texttt{simp} trades traceability for speed, so this book reaches for it only when a genuine technical obstacle makes the explicit alternative not worth the detour. A repeated \texttt{rw} chain or an independently-statable sub-goal is a lemma waiting to be named, but the converse is not a rule, splitting every proof into the smallest possible named lemmas on the mere \emph{possibility} of future reuse is the same mistake in the opposite direction; three similar lines are better than a premature abstraction, and the signal to extract is an actual repeated shape or an actual independently-statable sub-claim, not speculation.

